\documentclass[aspectratio=169,10pt]{beamer}
% ═══════════════════════════════════════════════════════════════════════════════
% SHARED PREAMBLE — Foundation Models Course (HIT)
% To be \input by each lecture file AFTER \documentclass
% ═══════════════════════════════════════════════════════════════════════════════

% ─────────────────────────────────────────────────────────────────────────────
% BEAMER THEME & BASIC SETUP
% ─────────────────────────────────────────────────────────────────────────────
\usetheme{Madrid}
\usecolortheme{default}

% Remove navigation symbols
\beamertemplatenavigationsymbolsempty

% ─────────────────────────────────────────────────────────────────────────────
% COLORS — HIT Branding
% ─────────────────────────────────────────────────────────────────────────────
\definecolor{hitblue}{RGB}{0,51,102}        % Primary: HIT Navy
\definecolor{hitgold}{RGB}{192,148,0}       % Accent: HIT Gold
\definecolor{hitlight}{RGB}{230,240,250}    % Light background

% Override Beamer palette
\setbeamercolor{primary}{fg=white,bg=hitblue}
\setbeamercolor{secondary}{fg=hitblue,bg=hitlight}
\setbeamercolor{structure}{fg=hitblue}
\setbeamercolor{alerted text}{fg=hitgold}

% ─────────────────────────────────────────────────────────────────────────────
% PACKAGES
% ─────────────────────────────────────────────────────────────────────────────
\usepackage{amsmath}
\usepackage{amssymb}
\usepackage{mathtools}
\usepackage{booktabs}
\usepackage{graphicx}
\usepackage{hyperref}
\usepackage{tikz}
\usepackage{pgfplots}
\usepackage{tcolorbox}
\usepackage{listings}
\usepackage{xcolor}
\usepackage{algorithm}
\usepackage{algpseudocode}

% ─────────────────────────────────────────────────────────────────────────────
% TikZ LIBRARIES
% ─────────────────────────────────────────────────────────────────────────────
\usetikzlibrary{shapes,arrows,positioning,calc,chains,scopes}
\pgfplotsset{compat=1.17}

% ─────────────────────────────────────────────────────────────────────────────
% CODE LISTINGS
% ─────────────────────────────────────────────────────────────────────────────
\lstset{
  basicstyle=\ttfamily\small,
  breaklines=true,
  commentstyle=\color{gray},
  keywordstyle=\color{hitblue},
  stringstyle=\color{hitgold},
  showstringspaces=false,
  backgroundcolor=\color{hitlight},
  frame=single,
  rulecolor=\color{hitblue}
}

% ─────────────────────────────────────────────────────────────────────────────
% TCOLORBOX STYLES
% ─────────────────────────────────────────────────────────────────────────────
\tcbset{
  hitbox/.style={
    colback=hitlight,
    colframe=hitblue,
    fonttitle=\bfseries,
    boxrule=1pt
  },
  alertbox/.style={
    colback=yellow!10,
    colframe=hitgold,
    fonttitle=\bfseries,
    boxrule=1.5pt
  }
}

% ─────────────────────────────────────────────────────────────────────────────
% CUSTOM COMMANDS
% ─────────────────────────────────────────────────────────────────────────────

% Placeholder for figures
\newcommand{\placeholder}[3]{
  \begin{center}
    \fbox{\begin{minipage}[c][#2][c]{#1}
      \centering\small\color{gray}[Figure: #3]
    \end{minipage}}
  \end{center}
}

% Section divider frame
\newcommand{\sectionframe}[2]{
  \begin{frame}[plain]
    \begin{tikzpicture}[remember picture,overlay]
      \fill[hitblue] (current page.south west) rectangle (current page.north east);
      \node[anchor=center, white, font=\Large\bfseries] at (current page.center) {
        \begin{tabular}{c}
          #1 \\[0.3cm]
          {\small\color{hitgold}#2}
        \end{tabular}
      };
    \end{tikzpicture}
  \end{frame}
}

% ─────────────────────────────────────────────────────────────────────────────
% LOAD CUSTOM MACROS
% ─────────────────────────────────────────────────────────────────────────────
% ═══════════════════════════════════════════════════════════════════════════════
% SHARED MACROS — Mathematical Notation for Foundation Models Course
% ═══════════════════════════════════════════════════════════════════════════════

% ─────────────────────────────────────────────────────────────────────────────
% ACTIVATION & LAYER FUNCTIONS
% ─────────────────────────────────────────────────────────────────────────────
\newcommand{\softmax}{\mathrm{softmax}}
\newcommand{\sigmoid}{\mathrm{sigmoid}}
\newcommand{\relu}{\mathrm{ReLU}}
\newcommand{\gelu}{\mathrm{GELU}}
\newcommand{\LayerNorm}{\mathrm{LayerNorm}}
\newcommand{\FFN}{\mathrm{FFN}}
\newcommand{\MHA}{\mathrm{MHA}}
\newcommand{\Attn}{\mathrm{Attention}}

% ─────────────────────────────────────────────────────────────────────────────
% ATTENTION COMPONENTS
% ─────────────────────────────────────────────────────────────────────────────
\newcommand{\query}{Q}
\newcommand{\key}{K}
\newcommand{\val}{V}
\newcommand{\dmodel}{d_{\text{model}}}
\newcommand{\dk}{d_k}
\newcommand{\nheads}{h}

% Scaled dot-product attention formula
\newcommand{\SDPA}{
  \Attn(\query, \key, \val) = \softmax\left(\frac{\query \key^T}{\sqrt{\dk}}\right) \val
}

\newcommand{\CrossAttn}{\text{CrossAttention}}

% ─────────────────────────────────────────────────────────────────────────────
% PROBABILITY & EXPECTATION
% ─────────────────────────────────────────────────────────────────────────────
\newcommand{\E}{\mathbb{E}}
\newcommand{\Prob}{\mathbb{P}}
\newcommand{\KL}{\mathrm{KL}}
\newcommand{\ELBO}{\mathrm{ELBO}}
\newcommand{\given}{\mid}

% ─────────────────────────────────────────────────────────────────────────────
% NORMS & OPERATIONS
% ─────────────────────────────────────────────────────────────────────────────
\newcommand{\norm}[1]{\left\lVert #1 \right\rVert}
\newcommand{\abs}[1]{\left| #1 \right|}
\newcommand{\inner}[2]{\langle #1, #2 \rangle}
\newcommand{\T}{^{\top}}

% ─────────────────────────────────────────────────────────────────────────────
% VECTORS & MATRICES
% ─────────────────────────────────────────────────────────────────────────────
\newcommand{\bvec}[1]{\mathbf{#1}}
\newcommand{\mat}[1]{\mathbf{#1}}
\newcommand{\iid}{\stackrel{\text{i.i.d.}}{\sim}}
\newcommand{\defeq}{\coloneq}

% ─────────────────────────────────────────────────────────────────────────────
% LANGUAGE MODELING
% ─────────────────────────────────────────────────────────────────────────────
\newcommand{\vocab}{\mathcal{V}}
\newcommand{\context}{\text{context}}
\newcommand{\token}{t}
\newcommand{\seq}{x}
\newcommand{\ptheta}{p_\theta}
\newcommand{\pref}{p_{\text{ref}}}

% ─────────────────────────────────────────────────────────────────────────────
% RLHF & DPO
% ─────────────────────────────────────────────────────────────────────────────
\newcommand{\piref}{\pi_{\text{ref}}}
\newcommand{\pitheta}{\pi_\theta}
\newcommand{\yw}{y_w}
\newcommand{\yl}{y_l}
\newcommand{\DPOloss}{\mathcal{L}_{\mathrm{DPO}}}

% ─────────────────────────────────────────────────────────────────────────────
% RETRIEVAL & RAG
% ─────────────────────────────────────────────────────────────────────────────
\newcommand{\docs}{\mathcal{D}}
\newcommand{\qvec}{q}
\newcommand{\dvec}{d}

% ─────────────────────────────────────────────────────────────────────────────
% AGENTS & REASONING
% ─────────────────────────────────────────────────────────────────────────────
\newcommand{\obs}{o}
\newcommand{\act}{a}
\newcommand{\thought}{t}
\newcommand{\tools}{\mathcal{T}}
\newcommand{\history}{h}
\newcommand{\concat}{\oplus}

% ─────────────────────────────────────────────────────────────────────────────
% SETS & LOGIC
% ─────────────────────────────────────────────────────────────────────────────
\newcommand{\R}{\mathbb{R}}
\newcommand{\N}{\mathbb{N}}
\newcommand{\Z}{\mathbb{Z}}

\endinput


% ─────────────────────────────────────────────────────────────────────────────
% TITLE & AUTHOR (can be overridden in individual lectures)
% ─────────────────────────────────────────────────────────────────────────────
\author[D. Lederman]{Dror Lederman}
\institute{Holon Institute of Technology (HIT) \\ Data Science Department}
\date{\today}

% ─────────────────────────────────────────────────────────────────────────────
% FOOTER & HEADER
% ─────────────────────────────────────────────────────────────────────────────
\setbeamertemplate{footline}{
  \leavevmode%
  \hbox{%
    \begin{beamercolorbox}[wd=0.33\paperwidth,ht=2.25ex,dp=1ex,center]{author in head/foot}%
      \usebeamerfont{author in head/foot}\insertshortauthor
    \end{beamercolorbox}%
    \begin{beamercolorbox}[wd=0.34\paperwidth,ht=2.25ex,dp=1ex,center]{title in head/foot}%
      \usebeamerfont{title in head/foot}\insertshorttitle
    \end{beamercolorbox}%
    \begin{beamercolorbox}[wd=0.33\paperwidth,ht=2.25ex,dp=1ex,right]{frame number}%
      \usebeamerfont{frame number}\insertframenumber{} / \inserttotalframenumber\hspace*{2ex}
    \end{beamercolorbox}%
  }%
  \vskip0pt%
}

\endinput


\title{Week 8: Reasoning and Planning in LLM Agents}
\subtitle{Foundation Models and Agentic AI}
\author{Lecturer Name}
\institute{Holon Institute of Technology (HIT)}
\date{Semester, 2025}

\begin{document}

\begin{frame}
  \titlepage
\end{frame}

\begin{frame}{Learning Objectives}
  After this lecture you will be able to:
  \begin{itemize}
    \item Explain \textbf{chain-of-thought (CoT)} prompting and its variants.
    \item Describe \textbf{Tree of Thoughts (ToT)} and search-based reasoning.
    \item Understand \textbf{Reflexion} as verbal reinforcement learning via self-reflection.
    \item Identify the \textbf{limits of LLM reasoning} and when it breaks down.
    \item Discuss \textbf{process reward models} and verification strategies.
    \item Apply \textbf{self-consistency} to improve reliability without architecture changes.
  \end{itemize}
\end{frame}

\section{Chain-of-Thought Prompting}
\sectionframe{Chain-of-Thought}{Teaching models to think step by step}

\begin{frame}{Chain-of-Thought Prompting}
  \textbf{Wei et al.\ (2022)}~\cite{wei2022chain}
  \vspace{0.4em}
  \begin{tcolorbox}[hitbox, title=CoT Prompt Example]
    \textbf{Standard:} Q: Roger has 5 tennis balls. He buys 2 more cans (3 balls each). How many does he have?  A: 11.\\[6pt]
    \textbf{CoT:} Q: Roger has 5 tennis balls... A: \textit{Roger starts with 5. He buys $2 \times 3 = 6$ more. $5 + 6 = 11$.} A: 11.
  \end{tcolorbox}
  \vspace{0.4em}
  \textbf{Key findings:}
  \begin{itemize}
    \item CoT emerges as a useful ability only at $\sim$100B parameter scale.
    \item Zero-shot CoT: append ``\textit{Let's think step by step.}'' — surprisingly effective.
    \item Gains largest on \textbf{multi-step arithmetic and commonsense reasoning}.
  \end{itemize}
  \note{Ask students why CoT helps. One hypothesis: the intermediate tokens serve as a scratchpad that expands the model's effective computation. The KV-cache of the reasoning acts as working memory.}
\end{frame}

\begin{frame}{Self-Consistency: Majority Vote over CoT Paths}
  \textbf{Wang et al.\ (2023)}~\cite{wang2023selfconsistency}
  \vspace{0.4em}
  \begin{tcolorbox}[hitbox, title=Self-Consistency]
    Sample $k$ diverse CoT reasoning paths from the model.
    Take the \textbf{majority vote} on the final answers.
    \[
      \hat{y} = \argmax_y \sum_{i=1}^k \mathbb{1}[y_i = y]
    \]
  \end{tcolorbox}
  \vspace{0.4em}
  \begin{itemize}
    \item Greedy decoding ($k=1$) is brittle; a single wrong reasoning step derails the answer.
    \item Self-consistency with $k=40$: +17 points on GSM8K over greedy CoT.
    \item Trade-off: $k\times$ inference cost.
    \item Works best when \emph{most paths lead to the right answer}, even via different routes.
  \end{itemize}
\end{frame}

\section{Tree of Thoughts}
\sectionframe{Tree of Thoughts}{Deliberate problem solving via search}

\begin{frame}{Tree of Thoughts (ToT)}
  \textbf{Yao et al.\ (2023)}~\cite{yao2023tree}
  \vspace{0.4em}
  \begin{columns}[T]
    \column{0.55\textwidth}
      \textbf{Key idea:}
      \begin{itemize}
        \item Decompose problems into \textbf{thoughts} (intermediate steps).
        \item At each step: generate $k$ candidate thoughts.
        \item \textbf{Evaluate} each thought (LLM-as-judge: ``sure/maybe/impossible'').
        \item \textbf{Search} over the tree: BFS or DFS.
      \end{itemize}
      \vspace{0.4em}
      \textbf{Results on Game of 24:}\\
      ToT: 74\% vs.\ CoT: 4\% (solving a card game requiring 4 arithmetic operations).
    \column{0.41\textwidth}
      \placeholder{0.95\textwidth}{4.5cm}{ToT: tree of thought nodes with evaluation scores}
  \end{columns}
  \note{The Game of 24 result is dramatic. The gap illustrates that CoT linear reasoning fundamentally fails on search problems.}
\end{frame}

\begin{frame}{ToT: Components}
  \begin{columns}[T]
    \column{0.52\textwidth}
      \textbf{Thought generation:}
      \begin{itemize}
        \item \emph{Sample}: draw $k$ thoughts i.i.d.
        \item \emph{Propose}: prompt for a list of $k$ distinct candidates.
      \end{itemize}
      \vspace{0.4em}
      \textbf{State evaluation:}
      \begin{itemize}
        \item \emph{Value}: score each thought independently with LLM.
        \item \emph{Vote}: LLM chooses the best among candidates.
      \end{itemize}
      \vspace{0.4em}
      \textbf{Search algorithm:}
      \begin{itemize}
        \item \emph{BFS}: expand best $b$ nodes at each depth.
        \item \emph{DFS}: explore one path to completion; backtrack on failure.
      \end{itemize}
    \column{0.44\textwidth}
      \begin{algorithm}[H]
        \caption{ToT BFS}
        \begin{algorithmic}[1]
          \State $S \gets \{x\}$
          \For{$t = 1$ to $T$}
            \State $S' \gets \bigcup_{s \in S} G(s, k)$
            \State $V \gets \text{Eval}(S')$
            \State $S \gets \text{TopB}(S', V, b)$
          \EndFor
          \Return $\text{Best}(S)$
        \end{algorithmic}
      \end{algorithm}
  \end{columns}
\end{frame}

\section{Reflexion}
\sectionframe{Reflexion}{Verbal reinforcement learning via self-reflection}

\begin{frame}{Reflexion: Learning from Failure}
  \textbf{Shinn et al.\ (2023)}~\cite{shinn2023reflexion}
  \vspace{0.4em}
  \begin{tcolorbox}[hitbox, title=Reflexion Loop]
    \begin{enumerate}
      \item Run agent; receive task feedback (success/failure/observation).
      \item \textbf{Reflect}: LLM writes a verbal self-evaluation in natural language.
      \item \textbf{Store} the reflection in an episodic memory buffer.
      \item \textbf{Re-attempt} the task with reflections in context.
    \end{enumerate}
  \end{tcolorbox}
  \vspace{0.3em}
  \textbf{Key analogy:} Reflexion is verbal RL — the ``gradient'' is a natural language critique, stored in memory rather than model weights.
  \vspace{0.3em}\\
  \textbf{Results:} +22\% on HumanEval coding, +20\% on AlfWorld decision-making.
\end{frame}

\begin{frame}{Reflexion: Example Trace}
  \begin{tcolorbox}[hitbox]
    \textbf{Attempt 1:} [code that fails test case 3] \\[4pt]
    \textbf{Reflection:} ``My solution assumed all input lists are non-empty, but test case 3 provides an empty list.
    In the next attempt, I should add a guard for empty input at the start of the function.''\\[4pt]
    \textbf{Attempt 2:} [code with empty-list guard — passes all tests] ✓
  \end{tcolorbox}
  \vspace{0.4em}
  \textbf{What's stored in memory:} The \emph{textual} reflection, not a gradient.\\
  \textbf{Limitation:} Memory grows with failed attempts; context length is eventually exhausted.
\end{frame}

\section{Limits of LLM Reasoning}
\sectionframe{Limits of LLM Reasoning}{When do LLMs fail to reason?}

\begin{frame}{Known Failure Modes}
  \begin{columns}[T]
    \column{0.52\textwidth}
      \textbf{Systematic failures:}
      \begin{itemize}
        \item \textbf{Reversal curse}: models trained on ``A is B'' struggle with ``B is A''.
        \item \textbf{Compositionality}: fail to reliably combine multiple rules.
        \item \textbf{Counting and arithmetic}: errors increase sharply with digit length.
        \item \textbf{Spatial reasoning}: 3D manipulation, maze navigation.
        \item \textbf{Planning horizon}: performance degrades with plan length.
        \item \textbf{Sycophancy}: agree with user pressure, even when wrong.
      \end{itemize}
    \column{0.44\textwidth}
      \textbf{Root causes:}
      \begin{itemize}
        \item Next-token prediction is not planning.
        \item No dedicated working memory outside context.
        \item Training distribution biases.
        \item RLHF alignment incentivizes agreement.
      \end{itemize}
      \placeholder{0.95\textwidth}{2.5cm}{GSM8K accuracy vs.\ number of digit operations}
  \end{columns}
\end{frame}

\begin{frame}{Process Reward Models (PRMs)}
  \begin{tcolorbox}[hitbox, title=Outcome vs.\ Process Reward]
    \textbf{Outcome Reward Model (ORM):} reward the final answer only.\\
    \textbf{Process Reward Model (PRM):} reward each intermediate \emph{reasoning step}.
  \end{tcolorbox}
  \vspace{0.4em}
  \textbf{Why PRMs?}
  \begin{itemize}
    \item ORM cannot distinguish a correct answer reached via wrong reasoning.
    \item PRM provides denser training signal.
    \item OpenAI PRM800K dataset: 800K step-level human labels on math problems.
    \item PRMs enable \textbf{Best-of-N}: generate $N$ solutions, pick highest step-score.
  \end{itemize}
  \vspace{0.3em}
  \textbf{Trade-off:} PRMs require expensive step-level annotation.
\end{frame}

\section{Paper Discussions}
\begin{frame}{Paper Discussion 1}
  \papercite{Chain-of-Thought Prompting Elicits Reasoning in Large Language Models}
    {Wei, Wang, Schuurmans et al.\ (Google Brain)}
    {NeurIPS 2022~\cite{wei2022chain}}
  \vspace{0.4em}
  \textbf{Key insight:} reasoning ability is a scale-emergent property.\\
  \textbf{Discussion:} If CoT only works at large scale, what does this imply for small models at the edge?
\end{frame}

\begin{frame}{Paper Discussion 2}
  \papercite{Reflexion: Language Agents with Verbal Reinforcement Learning}
    {Shinn, Cassano, Labash, Gopinath et al.}
    {NeurIPS 2023~\cite{shinn2023reflexion}}
  \vspace{0.4em}
  \textbf{Key contribution:} verbal critique stored in memory as a proxy for RL gradients.\\
  \textbf{Discussion:} Reflexion requires a reliable task feedback signal. For which tasks
  is binary success/failure feedback sufficient? For which tasks is it insufficient?
\end{frame}

\section{Lab / Demo}
\begin{frame}{Lab Idea: Comparative Reasoning Experiments}
  \begin{tcolorbox}[hitbox, title=Lab 8 — CoT vs.\ ToT vs.\ Reflexion]
    \textbf{Goal:} Compare reasoning strategies on a hard math/coding benchmark.\\[4pt]
    \textbf{Tools:} Python, OpenAI API, GSM8K or MATH dataset (50 problems)
    \begin{enumerate}
      \item Baseline: direct answer (no CoT).
      \item CoT: ``Let's think step by step.''
      \item Self-consistency: CoT × 5, majority vote.
      \item Reflexion: 3 retry attempts with verbal reflection.
      \item Measure accuracy and cost (API tokens) for each strategy.
    \end{enumerate}
    \textbf{Deliverable:} Results table + analysis of per-strategy failure modes.
  \end{tcolorbox}
\end{frame}

\section{Discussion \& Summary}
\begin{frame}{Discussion Questions}
  \begin{enumerate}
    \item Chain-of-thought works by having the model ``write out'' its reasoning. Does this mean the model is ``actually reasoning,'' or is it just generating tokens that look like reasoning? How would you test the difference?
    \item Tree of Thoughts requires an LLM to evaluate its own intermediate steps. When does the evaluator LLM fail? Can a model reliably evaluate steps harder than it can generate?
    \item Reflexion stores verbal self-critiques in the context window. What happens when the context fills up? How would you design a memory-efficient version of Reflexion?
    \item Process reward models require human step-level annotation of reasoning chains. Is this scalable? What are the alternatives?
  \end{enumerate}
\end{frame}

\begin{frame}{Summary}
  \begin{itemize}
    \item \textbf{CoT} prompting unlocks multi-step reasoning at scale; zero-shot CoT works with ``think step by step.''
    \item \textbf{Self-consistency} improves CoT via majority vote over diverse reasoning paths.
    \item \textbf{Tree of Thoughts} enables deliberate search; critical for planning and combinatorial tasks.
    \item \textbf{Reflexion} implements verbal RL: failures are verbalized and stored as memory.
    \item LLMs have systematic reasoning limits: reversal, compositionality, planning horizon.
    \item \textbf{PRMs} enable dense step-level supervision; used in OpenAI o1 and similar systems.
  \end{itemize}
  \vspace{0.3em}
  \textbf{Next week:} Memory architectures for agentic systems.
\end{frame}

\begin{frame}[allowframebreaks]{Suggested Reading}
  \nocite{wei2022chain, yao2023tree, shinn2023reflexion, wang2023selfconsistency}
  \bibliography{../../references}
\end{frame}

\end{document}
